%==========================================================
%  CHAPITRE 6 — FONCTIONS LOGARITHMIQUES
%==========================================================

\chapter{Fonctions Logarithmiques}

\begin{rappel}[Lien avec la fonction exponentielle]
  \begin{itemize}
    \item $\ln$ et $\exp$ sont \textbf{r\'{e}ciproques} l'une de l'autre :
          $e^{\ln x} = x$ et $\ln(e^x) = x$.
    \item $e^x > 0$ pour tout $x \in \R$.
    \item $\ln$ est d\'{e}finie uniquement sur $]0,+\infty[$.
  \end{itemize}
\end{rappel}

%----------------------------------------------------------
\section{Fonction Logarithme N\'{e}p\'{e}rien}
%----------------------------------------------------------

\subsection{D\'{e}finition}

\begin{definition}[Logarithme n\'{e}p\'{e}rien]
  La \textbf{fonction logarithme n\'{e}p\'{e}rien}, not\'{e}e $\ln$,
  est l'unique primitive de $x \mapsto \dfrac{1}{x}$ sur $]0,+\infty[$
  qui s'annule en $1$ :
  \[
    (\ln x)' = \frac{1}{x} \quad \text{sur } ]0,+\infty[
    \qquad \text{et} \qquad \ln(1) = 0
  \]
\end{definition}

\begin{propriete}[Cons\'{e}quences directes]
  \[
    \ln(1) = 0 \qquad \ln(e) = 1 \qquad
    \ln(x) = 0 \iff x = 1 \qquad
    \ln(x) > 0 \iff x > 1
  \]
\end{propriete}

%----------------------------------------------------------
\subsection{R\`{e}gles de calcul}
%----------------------------------------------------------

\begin{propriete}[Propri\'{e}t\'{e}s alg\'{e}briques]
  Pour tous $x > 0$, $y > 0$ et $r \in \Q$ :

  \bigskip
  \begin{center}
  \renewcommand{\arraystretch}{2}
  \begin{tabular}{|c|c|}
    \hline
    \textbf{\color{navy}Propri\'{e}t\'{e}}
      & \textbf{\color{navy}Formule} \\
    \hline\hline
    Produit
      & $\ln(xy) = \ln x + \ln y$ \\
    \hline
    Quotient
      & $\ln\!\left(\dfrac{x}{y}\right) = \ln x - \ln y$ \\
    \hline
    Inverse
      & $\ln\!\left(\dfrac{1}{x}\right) = -\ln x$ \\
    \hline
    Puissance
      & $\ln(x^r) = r\ln x$ \\
    \hline
    Racine carr\'{e}e
      & $\ln(\sqrt{x}) = \dfrac{1}{2}\ln x$ \\
    \hline
  \end{tabular}
  \end{center}
\end{propriete}

\begin{propriete}[Comparaison et \'{e}quations]
  Pour $x > 0$ et $y > 0$ :
  \[
    \ln x = \ln y \iff x = y \qquad
    \ln x < \ln y \iff x < y
  \]
  \[
    \ln x = y \iff x = e^y \qquad
    \ln x = r \iff x = e^r \quad (r \in \R)
  \]
\end{propriete}

\begin{attention}
  \textbf{Erreurs fr\'{e}quentes \`{a} \'{e}viter :}
  \[
    \ln(x + y) \neq \ln x + \ln y \qquad
    \ln(x \times y) \neq \ln x \times \ln y \qquad
    \frac{\ln x}{\ln y} \neq \ln\!\left(\frac{x}{y}\right)
  \]
\end{attention}

%----------------------------------------------------------
\subsection{Domaine de d\'{e}finition}
%----------------------------------------------------------

\begin{propriete}[Domaine des fonctions avec $\ln$]
  \bigskip
  \begin{center}
  \renewcommand{\arraystretch}{2}
  \begin{tabular}{|c|c|}
    \hline
    \textbf{\color{navy}Fonction}
      & \textbf{\color{navy}Domaine de d\'{e}finition} \\
    \hline\hline
    $f(x) = \ln x$
      & $D_f = \;]0,+\infty[$ \\
    \hline
    $f(x) = \ln\bigl(u(x)\bigr)$
      & $D_f = \{x \in D_u \mid u(x) > 0\}$ \\
    \hline
    $f(x) = \ln\bigl(u(x)^2\bigr)$
      & $D_f = \{x \in D_u \mid u(x) \neq 0\}$ \\
    \hline
    $f(x) = \ln\bigl|u(x)\bigr|$
      & $D_f = \{x \in D_u \mid u(x) \neq 0\}$ \\
    \hline
  \end{tabular}
  \end{center}
\end{propriete}

\begin{exemple}[R\'{e}solus : domaines de d\'{e}finition]
  \begin{itemize}
    \item $f(x) = \ln(2x - 3)$ :
    \quad $2x - 3 > 0 \Rightarrow x > \dfrac{3}{2}$
    \quad $\Rightarrow D_f = \left]\dfrac{3}{2},+\infty\right[$

    \medskip
    \item $f(x) = \ln(x^2 - 4)$ :
    \quad $x^2 - 4 > 0 \Rightarrow (x-2)(x+2) > 0$
    \quad $\Rightarrow D_f = \;]-\infty,-2[\;\cup\;]2,+\infty[$

    \medskip
    \item $f(x) = \ln(x^2 + 1)$ :
    \quad $x^2 + 1 > 0$ toujours vrai
    \quad $\Rightarrow D_f = \R$
  \end{itemize}
\end{exemple}

%----------------------------------------------------------
\subsection{Limites usuelles}
%----------------------------------------------------------

\begin{propriete}[Limites de r\'{e}f\'{e}rence]
  \bigskip
  \begin{center}
  \renewcommand{\arraystretch}{2}
  \begin{tabular}{|c|c|}
    \hline
    \textbf{\color{navy}Limite}
      & \textbf{\color{navy}R\'{e}sultat} \\
    \hline\hline
    $\displaystyle\lim_{x \to +\infty} \ln x$
      & $+\infty$ \\
    \hline
    $\displaystyle\lim_{x \to 0^+} \ln x$
      & $-\infty$ \\
    \hline
    $\displaystyle\lim_{x \to +\infty} \dfrac{\ln x}{x^n}$
      & $0$ \quad ($n \in \N^*$) \quad \textit{(croissance compar\'{e}e)} \\
    \hline
    $\displaystyle\lim_{x \to 0^+} x^n \ln x$
      & $0$ \quad ($n \in \N^*$) \quad \textit{(croissance compar\'{e}e)} \\
    \hline
    $\displaystyle\lim_{x \to 1} \dfrac{\ln x}{x - 1}$
      & $1$ \\
    \hline
    $\displaystyle\lim_{x \to 0} \dfrac{\ln(1+x)}{x}$
      & $1$ \\
    \hline
  \end{tabular}
  \end{center}
\end{propriete}

%----------------------------------------------------------
\subsection{D\'{e}rivabilit\'{e}}
%----------------------------------------------------------

\begin{propriete}[D\'{e}riv\'{e}e de $\ln x$ et de $\ln u(x)$]
  \begin{itemize}
    \item La fonction $\ln$ est d\'{e}rivable sur $]0,+\infty[$ et :
    \[
      (\ln x)' = \frac{1}{x}
    \]
    \item Si $u$ est d\'{e}rivable et $u > 0$ sur $I$, alors :
    \[
      \bigl(\ln u(x)\bigr)' = \frac{u'(x)}{u(x)}
    \]
  \end{itemize}
\end{propriete}

\begin{exemple}[R\'{e}solus : d\'{e}riv\'{e}es]
  \begin{itemize}
    \item $f(x) = \ln(3x+1)$ :
    \quad $u = 3x+1$, $u' = 3$
    \quad $\Rightarrow f'(x) = \dfrac{3}{3x+1}$

    \medskip
    \item $f(x) = \ln(x^2+1)$ :
    \quad $u = x^2+1$, $u' = 2x$
    \quad $\Rightarrow f'(x) = \dfrac{2x}{x^2+1}$

    \medskip
    \item $f(x) = x\ln x$ :
    \quad r\`{e}gle du produit :
    $f'(x) = 1 \cdot \ln x + x \cdot \dfrac{1}{x} = \ln x + 1$

    \medskip
    \item $f(x) = \ln\!\left(\dfrac{x+1}{x-1}\right) = \ln(x+1) - \ln(x-1)$ :
    \quad $f'(x) = \dfrac{1}{x+1} - \dfrac{1}{x-1} = \dfrac{-2}{x^2-1}$
  \end{itemize}
\end{exemple}

%----------------------------------------------------------
\subsection{Signe et repr\'{e}sentation graphique}
%----------------------------------------------------------

\begin{propriete}[Signe de $\ln x$]
  \bigskip
  \begin{center}
  \renewcommand{\arraystretch}{1.8}
  \begin{tabular}{|c|ccccc|}
    \hline
    $x$      & $0$ && $1$ && $+\infty$ \\
    \hline
    $\ln x$  & \multicolumn{2}{c}{$-\infty \quad -$} & $0$ & $+$ & $+\infty$ \\
    \hline
  \end{tabular}
  \end{center}

  \medskip
  $\ln x < 0$ pour $0 < x < 1$ \quad et \quad $\ln x > 0$ pour $x > 1$.
\end{propriete}

\begin{center}
\begin{tikzpicture}
  \begin{axis}[
    axis lines  = center,
    xlabel      = {$x$},
    ylabel      = {$y$},
    xmin = -0.5, xmax = 6,
    ymin = -3,   ymax = 3,
    xtick       = {0,1,2,3,4,5},
    ytick       = {-3,-2,-1,0,1,2,3},
    grid        = major,
    grid style  = {gray!20},
    width = 10cm, height = 8cm,
  ]
    % Courbe ln(x)
    \addplot[navy, very thick, domain=0.05:6, samples=200]
      {ln(x)};
    \addlegendentry{$f(x)=\ln x$}

    % Courbe e^x (réciproque)
    \addplot[forest, thick, dashed, domain=-2.5:1.8, samples=150]
      {exp(x)};
    \addlegendentry{$f^{-1}(x)=e^x$}

    % Bissectrice y=x
    \addplot[burgundy, thin, dashed, domain=-0.5:3]
      {x};
    \addlegendentry{$y=x$}

    % Asymptote x=0
    \addplot[gray, dotted, domain=-3:3]
      coordinates {(0,-3)(0,3)};

    % Points remarquables
    \addplot[mark=*, mark size=3pt, navy]
      coordinates {(1,0)(2.718,1)};
    \node[below right, navy, font=\small] at (axis cs:1,0) {$(1,0)$};
    \node[right, navy, font=\small] at (axis cs:2.718,1) {$(e,1)$};

    % Zone négative
    \node[burgundy, font=\small\bfseries] at (axis cs:0.4,-1.5) {$\ln x < 0$};
    % Zone positive
    \node[forest, font=\small\bfseries] at (axis cs:4,0.8) {$\ln x > 0$};
  \end{axis}
\end{tikzpicture}
\end{center}

%----------------------------------------------------------
\section{Fonction Logarithme de base $a$}
%----------------------------------------------------------

\begin{definition}[Logarithme de base $a$]
  Pour $a \in \R^*_+$ et $a \neq 1$,
  la \textbf{fonction logarithme de base $a$} est d\'{e}finie par :
  \[
    \forall x \in\, ]0,+\infty[, \qquad
    \log_a(x) = \frac{\ln x}{\ln a}
  \]
  \textbf{Cas particulier :} si $a = 10$, $\log_{10}$ est le
  \textbf{logarithme d\'{e}cimal}, not\'{e} $\log$.
\end{definition}

\begin{propriete}[R\`{e}gles de calcul]
  Pour $x > 0$, $y > 0$, $r \in \Q$ :
  \[
    \log_a(xy) = \log_a x + \log_a y \qquad
    \log_a\!\left(\frac{x}{y}\right) = \log_a x - \log_a y
  \]
  \[
    \log_a(x^r) = r\log_a x \qquad
    \log_a 1 = 0 \qquad
    \log_a a = 1
  \]
  \[
    \log_a x = \log_a y \iff x = y \qquad
    \log_a x = r \iff x = a^r
  \]
\end{propriete}

\begin{propriete}[Variations et limites selon $a$]
  \bigskip
  \begin{center}
  \renewcommand{\arraystretch}{2}
  \begin{tabular}{|p{6cm}|p{6cm}|}
    \hline
    \textbf{\color{burgundy}$0 < a < 1$ \quad (d\'{e}croissante)}
      & \textbf{\color{forest}$a > 1$ \quad (croissante)} \\
    \hline\hline
    $\log_a x > \log_a y \iff x < y$
      & $\log_a x > \log_a y \iff x > y$ \\
    \hline
    $\displaystyle\lim_{x\to+\infty} \log_a x = -\infty$
      & $\displaystyle\lim_{x\to+\infty} \log_a x = +\infty$ \\
    \hline
    $\displaystyle\lim_{x\to 0^+} \log_a x = +\infty$
      & $\displaystyle\lim_{x\to 0^+} \log_a x = -\infty$ \\
    \hline
    \multicolumn{2}{|c|}{%
      $\displaystyle\lim_{x\to 0} \dfrac{a^x - 1}{x} = \ln a$
      \qquad
      $\bigl(\log_a x\bigr)' = \dfrac{1}{x \ln a}$
    } \\
    \hline
  \end{tabular}
  \end{center}
\end{propriete}

\begin{center}
\begin{tikzpicture}
  \begin{axis}[
    axis lines  = center,
    xlabel      = {$x$},
    ylabel      = {$y$},
    xmin = -0.3, xmax = 6,
    ymin = -3,   ymax = 3,
    xtick       = {0,1,2,3,4,5},
    ytick       = {-3,-2,-1,0,1,2},
    grid        = major,
    grid style  = {gray!20},
    width = 10cm, height = 8cm,
  ]
    % log base 2 (a>1, croissante)
    \addplot[navy, very thick, domain=0.05:6, samples=200]
      {ln(x)/ln(2)};
    \addlegendentry{$\log_2 x$ \; ($a=2>1$)}

    % log base 1/2 (0<a<1, décroissante)
    \addplot[burgundy, very thick, domain=0.05:6, samples=200]
      {ln(x)/ln(0.5)};
    \addlegendentry{$\log_{1/2} x$ \; ($a=1/2<1$)}

    % ln (référence)
    \addplot[forest, thick, dashed, domain=0.05:6, samples=200]
      {ln(x)};
    \addlegendentry{$\ln x$ \; (r\'{e}f\'{e}rence, $a=e$)}

    % Point commun (1,0)
    \addplot[mark=*, mark size=4pt, black]
      coordinates {(1,0)};
    \node[above right, black, font=\small\bfseries]
      at (axis cs:1,0) {$(1,0)$};
  \end{axis}
\end{tikzpicture}
\end{center}

\begin{exemple}[R\'{e}solus : logarithme de base $a$]
  \begin{itemize}
    \item $\log_2(32) = \log_2(2^5) = 5$
    \medskip
    \item $\log_3(x) = 2 \Rightarrow x = 3^2 = 9$
    \medskip
    \item $\log_{10}(1000) = \log_{10}(10^3) = 3$
    \medskip
    \item R\'{e}soudre $\log_2 x > 3$ :
    \quad $x > 2^3 = 8$ \quad (car $a = 2 > 1$, $\log_2$ croissante)
  \end{itemize}
\end{exemple}

\newpage
%----------------------------------------------------------
\section*{Exercices du chapitre}
\addcontentsline{toc}{section}{Exercices : Fonctions logarithmiques}
%----------------------------------------------------------

\begin{exercice}[1 : Calculs avec $\ln$]
  Simplifier les expressions suivantes :
  \begin{multicols}{2}
  \begin{enumerate}
    \item $\ln(e^5)$
    \item $\ln\!\left(\dfrac{1}{e^3}\right)$
    \item $e^{2\ln 3}$
    \item $\ln(\sqrt{e})$
    \item $\ln(2) + \ln(3) - \ln(6)$
    \item $3\ln(2) - \ln(4)$
  \end{enumerate}
  \end{multicols}
\end{exercice}

\begin{exercice}[2 : Domaine de d\'{e}finition]
  D\'{e}terminer le domaine de d\'{e}finition de :
  \begin{enumerate}
    \item $f(x) = \ln(5x - 2)$
    \item $f(x) = \ln(x^2 - 9)$
    \item $f(x) = \ln(x^2 + 2x + 1)$
    \item $f(x) = \ln(x) + \ln(1-x)$
  \end{enumerate}
\end{exercice}

\begin{exercice}[3 : D\'{e}riv\'{e}es]
  Calculer la d\'{e}riv\'{e}e des fonctions suivantes :
  \begin{multicols}{2}
  \begin{enumerate}
    \item $f(x) = \ln(x^2 + 3)$
    \item $f(x) = x^2 \ln x$
    \item $f(x) = \dfrac{\ln x}{x}$
    \item $f(x) = \ln\!\left(\dfrac{x+2}{x-1}\right)$
    \item $f(x) = \ln(\cos x)$
    \item $f(x) = \ln(x + \sqrt{x^2+1})$
  \end{enumerate}
  \end{multicols}
\end{exercice}

\begin{exercice}[4 : \'{E}quations et in\'{e}quations]
  R\'{e}soudre dans $\R$ :
  \begin{enumerate}
    \item $\ln(2x+1) = \ln(x+3)$
    \item $\ln(x^2) = \ln(4)$
    \item $\ln(x) + \ln(x-2) = \ln(3)$
    \item $\ln^2(x) - 3\ln(x) + 2 = 0$
          \quad (poser $X = \ln x$)
    \item $\ln(x) > \ln(3x - 2)$
  \end{enumerate}
\end{exercice}

\begin{exercice}[5 : Limites]
  Calculer les limites suivantes :
  \begin{multicols}{2}
  \begin{enumerate}
    \item $\displaystyle\lim_{x\to+\infty} \frac{\ln x}{x}$
    \item $\displaystyle\lim_{x\to 0^+} x \ln x$
    \item $\displaystyle\lim_{x\to+\infty} \frac{\ln x}{\sqrt{x}}$
    \item $\displaystyle\lim_{x\to 0} \frac{\ln(1+3x)}{x}$
  \end{enumerate}
  \end{multicols}
\end{exercice}

\begin{exercice}[6 : \'{E}tude compl\`{e}te]
  Soit $f(x) = x - \ln x$ sur $]0,+\infty[$.
  \begin{enumerate}
    \item Calculer les limites de $f$ en $0^+$ et en $+\infty$.
    \item Calculer $f'(x)$ et dresser le tableau de variations.
    \item D\'{e}terminer le minimum de $f$.
    \item Montrer que $\ln x \leq x - 1$ pour tout $x > 0$.
  \end{enumerate}
\end{exercice}